%% kapa: exam=1 round=37 year=2026 period=1 form=A subject=real_estate subject_ko=부동산학원론 seq=37 number=117
%% source: Q-Net 감정평가사 기출문제 (cst00309) artlSeq=5258924 file=raw/1008/5258924/2026년 제37회 감정평가사 1차 1교시.pdf
%% conversion: kordoc 4.24.51 -> tools/gen1.py (auto); review=proofread-sonnet; flags=none
\begin{question}{117}
감정평가 실무기준상 건물의 감정평가방법으로 옳지 않은 것은?
\begin{choices}
  \item 재조달원가를 간접법으로 구하는 경우에는 건물신축단가표와 비교하거나 비슷한 건물의 신축원가 사례를 조사한 후 사정보정 및 시점수정 등을 하여 대상 건물의 재조달원가를 산정할 수 있다.
  \item 건물의 일부가 도시ㆍ군계획시설에 저촉되어 저촉되지 않은 잔여부분이 건물로서 효용가치가 없는 경우에는 감정평가에서 제외하고, 잔여부분만으로도 독립건물로서의 가치가 있다고 인정되는 경우에는 그 잔여부분의 벽체나 기둥 등의 보수에 드는 비용 등을 고려하여 감정평가한다.
  \item 건물의 일반적인 효용을 위한 전기설비, 승강기설비, 소화전설비 등 부대설비는 건물에 포함하여 감정평가한다. 다만, 특수한 목적의 경우에는 구분하여 감정평가할 수 있다.
  \item 기존 건물 상층부에 증축한 건물 증축부분의 장래보존연수는 기존 건물의 장래보존연수 범위에서 적용하여 감가수정한다.
  \item 실지조사 시 의뢰되지 않은 공부상 미등재 건물이 있는 경우에는 의뢰인에게 감정평가 포함 여부를 확인하여 실측면적을 기준으로 감정평가할 수 있다.
\end{choices}
\end{question}
\answer{2}
